% !TeX root = ../../Book4.tex
% 纲要标题：### 3.2 Yoneda 引理
% 纲要位置：计划纲要.md，第 3165 行。

\section{Yoneda 引理}
\label{b4:sec:0302}

一个对象接受的全部映射，是否已经决定这个对象及其间的态射？Yoneda 引理通过在恒等态射处求值回答这个问题，并给出所需自然性。

% 正文按纲要写入；各节的基础结果在本章证明。

% 纲要内容（仅作写作计划，尚非教材正文）：
%
% * Yoneda 嵌入；小范畴中的全忠实函子与局部小范畴的 Hom 双射
% * 引理的证明及自然性
% * 用 Hom 函子探测对象与态射
%

\subsection{Yoneda 嵌入}
\label{b4:subsec:0302:01}

对局部小范畴 \(\mathcal C\) 中的对象 \(A\)，函子 \(h_A=\operatorname{Hom}_{\mathcal C}(-,A)\) 记录从每个对象到 \(A\) 的态射。若 \(a:A\rightarrow B\)，后复合给出自然变换
\[
h_a:h_A\Longrightarrow h_B,
\qquad (h_a)_X(f)=a\circ f.
\]
自然性来自 \(a\circ(f\circ u)=(a\circ f)\circ u\)，其中 \(u:X'\rightarrow X\)。又有 \(h_{1_A}=1_{h_A}\) 及 \(h_{ba}=h_b\circ h_a\)。当 \(\mathcal C\) 小时，因此得到函子
\[
\mathsf y:\mathcal C\longrightarrow\operatorname{Fun}(\mathcal C^{\mathrm{op}},\mathbf{Set}),
\qquad A\longmapsto h_A,\quad a\longmapsto h_a.
\]
它称为\term[yoneda-qianru]{Yoneda 嵌入}[Yoneda embedding]。名称中的“嵌入”将在下面证明为全忠实性；一般不能把它理解成对象集合上的字面包含。
\symidx{\(\mathsf y\)：Yoneda 嵌入}

\ProseParagraphBreak
这里选用 \(h_A\)，是因为 \(a:A\rightarrow B\) 使 \(h_A\rightarrow h_B\)，在表示对象变量上同向变化。若改用 \(h^A=\operatorname{Hom}_{\mathcal C}(A,-)\)，预合成 \(a\) 得到的却是 \(h^B\rightarrow h^A\)，因此在小范畴情形得到的是 \(\mathcal C^{\mathrm{op}}\rightarrow\operatorname{Fun}(\mathcal C,\mathbf{Set})\)。

\ProseParagraphBreak
若 \(\mathcal C\) 仅局部小，下文按分量处理 \(h_A\) 以及涉及它的自然变换，不预设整个集合值函子范畴都在\chapref{b4:ch:01}的大小范围内。引理将证明，以 \(h_A\) 为源、\(F\) 为目标的自然变换由集合 \(F(A)\) 参数化，因而可表函子之间的自然变换构成集合。这些大小约定保证后面所写的 \(\operatorname{Nat}(h_A,F)\) 确实是集合；双射本身的构造则由恒等态射与自然性给出。对于仅局部小的范畴，本书所说的“Yoneda 全忠实性”指下文给出的 Hom 集与自然变换集合之间的双射。

\ProseParagraphBreak
在集合范畴，\(h_A(X)\) 是所有 \(X\) 指标的 \(A\) 中元素族；改变指标的映射 \(X'\rightarrow X\)，就是把元素族拉回。一个映射 \(a:A\rightarrow B\) 对每个族逐项应用 \(a\)，于是自动产生相容的族变换。对于一般范畴，态射 \(X\rightarrow A\) 取代了元素族。关键问题是：每一种对所有 \(X\) 都相容的变换，是否都这样得到？

\subsection{Yoneda 引理的证明与自然性}
\label{b4:subsec:0302:02}

对源、目标相同的两个函子 \(H,F\)，记 \(\operatorname{Nat}(H,F)\) 为从 \(H\) 到 \(F\) 的自然变换类。\term[yoneda-yinli]{Yoneda 引理}[Yoneda lemma]说明，以可表函子为源时，这个类是集合，而且只需计算一个分量中的一个值。

\begin{theorem}[Yoneda 引理]\label{b4:thm:yoneda}
设 \(\mathcal C\) 为局部小范畴，\(F:\mathcal C^{\mathrm{op}}\rightarrow\mathbf{Set}\) 为函子，\(A\in\mathcal C\)，记 \(h_A=\operatorname{Hom}_{\mathcal C}(-,A)\)。求值映射
\[
\Phi_{A,F}:\operatorname{Nat}(h_A,F)\longrightarrow F(A),
\qquad \alpha\longmapsto\alpha_A(1_A)
\]
是双射，因而自然变换类 \(\operatorname{Nat}(h_A,F)\) 是集合。其逆把 \(u\in F(A)\) 送到自然变换 \(\alpha^u\)，其中
\[
\alpha^u_X(f)=F(f)(u)\qquad(f:X\rightarrow A).
\]
这一双射在 \(A\) 与 \(F\) 上自然：设 \(a:A\rightarrow B\) 为 \(\mathcal C\) 中的态射，\(h_B=\operatorname{Hom}_{\mathcal C}(-,B)\)，\(h_a:h_A\Rightarrow h_B\) 为后复合 \(a\) 得到的自然变换。改变表示对象时，自然变换一边预合成 \(h_a\)，函子值一边通过 \(F(a)\) 变化。改变函子时，自然变换一边后复合，函子值一边作用相应的 \(A\) 分量。具体地，若 \(G:\mathcal C^{\mathrm{op}}\rightarrow\mathbf{Set}\) 为函子，\(\alpha:h_B\Rightarrow F\)、\(\beta:F\Rightarrow G\) 为自然变换，则
\[
\begin{aligned}
\Phi_{A,F}(\alpha\circ h_a)&=F(a)\bigl(\Phi_{B,F}(\alpha)\bigr),\\
\Phi_{A,G}(\beta\circ\gamma)&=\beta_A\bigl(\Phi_{A,F}(\gamma)\bigr)
\qquad(\gamma:h_A\Rightarrow F).
\end{aligned}
\]
对协变函子 \(F:\mathcal C\rightarrow\mathbf{Set}\)，记 \(h^A=\operatorname{Hom}_{\mathcal C}(A,-)\)，相应结论为
\(\operatorname{Nat}(h^A,F)\cong F(A)\)，逆映射仍由 \(f:A\rightarrow X\mapsto F(f)(u)\) 给出。
\end{theorem}
\begin{proof}
给定 \(u\in F(A)\)，首先验证公式确实定义自然变换。对 \(v:X'\rightarrow X\) 及 \(f:X\rightarrow A\)，反变函子公理给出
\[
\alpha^u_{X'}(fv)=F(fv)(u)=F(v)F(f)(u)
=F(v)\bigl(\alpha^u_X(f)\bigr).
\]
所以 \(\alpha^u\) 自然，且 \(\alpha^u_A(1_A)=u\)。反过来，若 \(\alpha:h_A\Rightarrow F\) 自然，任意 \(f:X\rightarrow A\) 都满足 \(h_A(f)(1_A)=1_A\circ f=f\)。所以所有这样的 \(f\) 都由同一个 \(1_A\) 经前复合得到。将自然性用于这条箭头和这个元素，得到
\[
\alpha_X(f)=F(f)\bigl(\alpha_A(1_A)\bigr).
\]
\[
\begin{tikzcd}[column sep=large,row sep=large]
h_A(A)\arrow[r,"\alpha_A"]\arrow[d,"h_A(f)"']&F(A)\arrow[d,"F(f)"]\\
h_A(X)\arrow[r,"\alpha_X"']&F(X).
\end{tikzcd}
\]
这说明从 \(\alpha_A(1_A)\) 按所给公式重建的自然变换就是 \(\alpha\)，故两个映射互逆。所有自然变换恰为由 \(u\in F(A)\) 构造出的 \(\alpha^u\)，所以它们构成一个由集合 \(F(A)\) 参数化的集合。

现在检查两个参数的自然性。对 \(a:A\rightarrow B\) 及 \(\alpha:h_B\Rightarrow F\)，有
\[
(\alpha\circ h_a)_A(1_A)=\alpha_A(a)
=F(a)\bigl(\alpha_B(1_B)\bigr).
\]
由 \(\Phi\) 在恒等态射处求值的定义，这正是下列方块在任意 \(\alpha\) 上的交换条件：
\[
\begin{tikzcd}[column sep=large,row sep=large]
\operatorname{Nat}(h_B,F)\arrow[r,"\Phi_{B,F}"]\arrow[d,"{(-)\circ h_a}"']&F(B)\arrow[d,"F(a)"]\\
\operatorname{Nat}(h_A,F)\arrow[r,"\Phi_{A,F}"']&F(A).
\end{tikzcd}
\]
对 \(\beta:F\Rightarrow G\) 及 \(\gamma:h_A\Rightarrow F\)，则
\[
(\beta\circ\gamma)_A(1_A)
=\beta_A\bigl(\gamma_A(1_A)\bigr).
\]
这同样说明后复合 \(\beta\) 与在函子值上作用 \(\beta_A\) 相容，即
\[
\begin{tikzcd}[column sep=large,row sep=large]
\operatorname{Nat}(h_A,F)\arrow[r,"\Phi_{A,F}"]\arrow[d,"{\beta\circ(-)}"']&F(A)\arrow[d,"\beta_A"]\\
\operatorname{Nat}(h_A,G)\arrow[r,"\Phi_{A,G}"']&G(A).
\end{tikzcd}
\]
协变形式对相反范畴 \(\mathcal C^{\mathrm{op}}\) 应用刚证的结论即可；在原范畴中，所用态射方向恰变成 \(A\rightarrow X\)。
\end{proof}

自然性迫使一个变换在所有态射上的值由它在表示对象的恒等态射上的值决定，函子公理则保证得到的公式在所有对象间相容。这就是前一节泛元素判据的公式为何适用于每个自然变换，而不只适用于自然同构。

\begin{corollary}\label{b4:cor:yoneda-embedding}
设 \(\mathcal C\) 为局部小范畴。对任意 \(A,B\in\mathcal C\)，记 \(h_A=\operatorname{Hom}_{\mathcal C}(-,A)\)、\(h_B=\operatorname{Hom}_{\mathcal C}(-,B)\)，以 \(h_a:h_A\Rightarrow h_B\) 表示后复合态射 \(a:A\rightarrow B\) 得到的自然变换。映射
\[
\operatorname{Hom}_{\mathcal C}(A,B)\longrightarrow\operatorname{Nat}(h_A,h_B),
\qquad a\longmapsto h_a
\]
是双射。若 \(\mathcal C\) 进一步为小范畴，则 Yoneda 嵌入
\[
\mathsf y:\mathcal C\longrightarrow\operatorname{Fun}(\mathcal C^{\mathrm{op}},\mathbf{Set}),
\qquad A\longmapsto h_A,\quad a\longmapsto h_a
\]
为全忠实函子。对于局部小范畴 \(\mathcal C\)，若 \(h_A\) 与 \(h_B\) 自然同构，则 \(A\cong B\)。
\end{corollary}
\begin{proof}
在 Yoneda 引理中取 \(F=h_B\)，右边为 \(h_B(A)=\operatorname{Hom}_{\mathcal C}(A,B)\)，逆对应正是后复合。因此态射映射为双射。当 \(\mathcal C\) 小时，函子 \(\mathsf y\) 由这些映射给出，因而全且忠实。若自然同构 \(\alpha:h_A\xrightarrow{\sim}h_B\) 及其逆分别来自 \(a:A\rightarrow B\)、\(b:B\rightarrow A\)，则
\[
h_{ba}=\alpha^{-1}\alpha=1_{h_A}=h_{1_A},\qquad
h_{ab}=1_{h_B}=h_{1_B}.
\]
上面双射的单射性给出 \(ba=1_A\)、\(ab=1_B\)，故 \(a\) 是同构。
\end{proof}

这一嵌入的标准表述可参阅 Weibel《An Introduction to Homological Algebra》\cite[Example A.3.4]{Weibel1994HomologicalAlgebra}。对于协变 Hom 函子，态射 \(s:B\rightarrow A\) 对应自然变换 \(h^A\Rightarrow h^B\)，其分量为
\[
\operatorname{Hom}_{\mathcal C}(A,X)\longrightarrow
\operatorname{Hom}_{\mathcal C}(B,X),\qquad f\longmapsto f\circ s.
\]
因此 \(\operatorname{Nat}(h^A,h^B)\cong\operatorname{Hom}_{\mathcal C}(B,A)\)。固定源的 Hom 函子在表示对象上反向变化，正是由于这里要预合成 \(s\)。

\subsection{用 Hom 函子探测对象与态射}
\label{b4:subsec:0302:03}

\begin{proposition}\label{b4:prop:hom-detects-morphisms}
设 \(\mathcal C\) 为局部小范畴，\(f:A\rightarrow B\) 为态射。则：\(f\) 为单态射，当且仅当每个 \(X\in\mathcal C\) 的后复合映射 \(\operatorname{Hom}_{\mathcal C}(X,A)\rightarrow\operatorname{Hom}_{\mathcal C}(X,B)\) 都单射；\(f\) 为满态射，当且仅当每个 \(X\) 的前复合映射 \(\operatorname{Hom}_{\mathcal C}(B,X)\rightarrow\operatorname{Hom}_{\mathcal C}(A,X)\) 都单射。此外，\(f\) 为同构，当且仅当上述后复合映射对每个 \(X\) 都双射。
\end{proposition}
\begin{proof}
前两项分别把单态射与满态射的消去条件按 Hom 集重写，没有增加条件。若 \(f\) 是同构，后复合 \(f^{-1}\) 给出每个映射的逆。反过来，取 \(X=B\)，该 Hom 集映射的满射性给出 \(g:B\rightarrow A\) 使 \(fg=1_B\)。再取 \(X=A\)，由于 \(f(gf)=f=f1_A\)，相应 Hom 集映射的单射性给出 \(gf=1_A\)。所以 \(f\) 是同构。
\end{proof}

这里必须保留“对每个 \(X\)”或另证明所选对象足以检验。在群范畴中，从平凡群到任意群都只有一个同态，仅用这个对象不能区分群。对左 \(R\) 模则可以用 \(R\) 检验同构，因为 \(\operatorname{Hom}_R(R,M)\cong U(M)\)，而双射的模同态具有线性的逆。这一特殊结论依赖具体范畴，不能当作所有范畴都有的底集合判据。

\begin{example}[交换群中的自然一元运算]\label{b4:exm:yoneda-abelian-operations}
设 \(U:\mathbf{Ab}\rightarrow\mathbf{Set}\) 为忘却函子。对每个交换群 \(A\)，给一个集合映射 \(t_A:U(A)\rightarrow U(A)\)，并要求对每个群同态 \(f:A\rightarrow B\) 有 \(f(t_A(a))=t_B(f(a))\)。则存在唯一整数 \(n\)，使所有 \(A\) 和 \(a\in A\) 都满足 \(t_A(a)=na\)。

\ProseParagraphBreak
事实上，取 \(n=t_{\mathbb Z}(1)\)。对唯一满足 \(f_a(1)=a\) 的同态 \(f_a:\mathbb Z\rightarrow A\)，自然性给出 \(t_A(a)=f_a(n)=na\)。反过来，乘固定整数与所有群同态交换。这里从未事先要求 \(t_A\) 为群同态；这一性质被对所有群的自然性迫出。用 Yoneda 语言，这是 \(U\cong h^{\mathbb Z}\) 与 \(\operatorname{Nat}(h^{\mathbb Z},U)\cong U(\mathbb Z)\) 的具体计算。
\end{example}

若只研究一个群 \(A\)，它的任意集合自映射都可以被指定；把这些映射要求为对全部群同态自然，选择就缩减为一个整数。Yoneda 引理的用途正在于把这种同时涉及许多对象的相容条件，转化成一个明确的参数集合。
